% temporal-programming.tex
% Temporal Programming: A Cell-Based Paradigm for
% Oscillator-Relative Distributed Computation
%
% Compile with:
%   pdflatex temporal-programming
%   bibtex temporal-programming
%   pdflatex temporal-programming
%   pdflatex temporal-programming

\documentclass[twocolumn,10pt]{article}

%-----------------------------------------------------------------------
% Packages
%-----------------------------------------------------------------------
\usepackage[margin=1in,columnsep=0.25in]{geometry}
\usepackage{amsmath,amssymb,amsthm}
\usepackage{mathtools}
\usepackage{algorithm}
\usepackage{algpseudocode}
\usepackage{listings}
\usepackage{booktabs}
\usepackage{graphicx}
\usepackage[colorlinks=true,linkcolor=blue,citecolor=blue,urlcolor=blue]{hyperref}
\usepackage{enumitem}
\usepackage[expansion=false]{microtype}
\usepackage{xcolor}
\usepackage{fancyhdr}
\usepackage{abstract}
\usepackage{titlesec}
\usepackage{caption}
\usepackage{subcaption}
\usepackage{cite}
\usepackage{url}
\usepackage{stmaryrd}   % \llbracket, \rrbracket for denotational brackets

%-----------------------------------------------------------------------
% Page style
%-----------------------------------------------------------------------
\pagestyle{fancy}
\fancyhf{}
\fancyhead[L]{\small\textit{Temporal Programming: Cell-Based Oscillator-Relative Computation}}
\fancyhead[R]{\small\thepage}
\renewcommand{\headrulewidth}{0.4pt}

%-----------------------------------------------------------------------
% Theorem environments
%-----------------------------------------------------------------------
\newtheorem{theorem}{Theorem}[section]
\newtheorem{lemma}[theorem]{Lemma}
\newtheorem{corollary}[theorem]{Corollary}
\newtheorem{proposition}[theorem]{Proposition}
\newtheorem{definition}[theorem]{Definition}
\newtheorem{remark}[theorem]{Remark}
\newtheorem{example}[theorem]{Example}

%-----------------------------------------------------------------------
% Listings / Tempus syntax
%-----------------------------------------------------------------------
\definecolor{tempuskw}{RGB}{0,0,180}
\definecolor{tempuscomment}{RGB}{60,128,49}
\definecolor{tempusstring}{RGB}{163,21,21}
\definecolor{tempusnumber}{RGB}{9,134,88}
\definecolor{codebg}{RGB}{248,248,248}

\lstdefinelanguage{Tempus}{
  morekeywords={cell,bounds,action,sync,at,compose,channels,into,when,do,
                ensemble,sync_threshold,let,in,if,then,else,emit,wait,
                period,freq,channel,fire,return,begin,end,true,false,
                COMPILE,EXECUTE,idle,assign},
  sensitive=true,
  morecomment=[l]{--},
  morecomment=[s]{(*}{*)},
  morestring=[b]",
  keywordstyle=\color{tempuskw}\bfseries,
  commentstyle=\color{tempuscomment}\itshape,
  stringstyle=\color{tempusstring},
  numberstyle=\tiny\color{tempusnumber},
  basicstyle=\footnotesize\ttfamily,
  backgroundcolor=\color{codebg},
  frame=single,
  framesep=4pt,
  xleftmargin=8pt,
  xrightmargin=4pt,
  breaklines=true,
  showstringspaces=false,
  tabsize=2,
  captionpos=b
}

\lstset{language=Tempus}

%-----------------------------------------------------------------------
% Section formatting
%-----------------------------------------------------------------------
\titleformat{\section}{\normalfont\large\bfseries}{\thesection}{1em}{}
\titleformat{\subsection}{\normalfont\normalsize\bfseries}{\thesubsection}{1em}{}
\titleformat{\subsubsection}{\normalfont\normalsize\itshape}{\thesubsubsection}{1em}{}

%-----------------------------------------------------------------------
% Math macros
%-----------------------------------------------------------------------
\newcommand{\DP}{\Delta P}          % timing deviation
\newcommand{\Tref}{T_{\mathrm{ref}}}
\newcommand{\trec}{t_{\mathrm{rec}}}
\newcommand{\Mcyc}{M}               % cycle count
\newcommand{\fref}{f_{\mathrm{ref}}}
\newcommand{\calC}{\mathcal{C}}     % cell
\newcommand{\calD}{\mathcal{D}}     % DP-space
\newcommand{\calA}{\mathcal{A}}     % action set
\newcommand{\calT}{\mathcal{T}}     % trajectory
\newcommand{\calP}{\mathcal{P}}     % partition
\newcommand{\calE}{\mathcal{E}}     % execution environment
\newcommand{\calS}{\mathcal{S}}     % state space
\newcommand{\llb}{\llbracket}
\newcommand{\rrb}{\rrbracket}
\newcommand{\dcomp}{\;\|\;}         % parallel composition
\DeclareMathOperator{\Leb}{Leb}
\DeclareMathOperator{\cell}{cell}
\DeclareMathOperator{\traj}{traj}
\DeclareMathOperator{\Act}{Act}
\DeclareMathOperator{\Cells}{Cells}
\DeclareMathOperator{\Phase}{Phase}

%-----------------------------------------------------------------------
% Title
%-----------------------------------------------------------------------
\title{%
  \LARGE\bfseries Temporal Programming:\\
  A Cell-Based Paradigm for Oscillator-Relative\\
  Distributed Computation%
}

\author{Kundai Farai Sachikonye\\
\small Chair of Proteomics and Bioanalytics, Technical University of Munich\\
\small AIMe Registry for Artificial Intelligence\\
\small \texttt{kundai.sachikonye@wzw.tum.de}}

\date{May 2026}

%=======================================================================
\begin{document}
%=======================================================================

\twocolumn[{%
\maketitle
\begin{@twocolumnfalse}

%-----------------------------------------------------------------------
\begin{abstract}
We introduce \emph{temporal programming}, a new programming paradigm in
which every inter-node signal is a bare clock pulse carrying no semantic
payload. The program datum is the \emph{timing deviation}
$\DP(k) = \Tref(k) - \trec(k)$: the signed difference between the
oscillator-predicted arrival time of the $k$-th pulse and its actual
reception time. Computation proceeds by classifying the sequence of
$n$ timing deviations into one of a finite set of \emph{timing cells}
--- measurable subsets of $\DP$-space with positive Lebesgue measure ---
and dispatching the cell-associated action. We prove the
\emph{Composition Inflation Theorem}: for $n$ timing events distributed
over $d$ active channels, the number of categorically distinguishable
labeled trajectories is $T(n,d) = d\,(1+d)^{n-1}$, growing
combinatorially in both arguments. We demonstrate that the paradigm
enforces \emph{structural incorruptibility}: because no payload parser
exists, arbitrary code injection, buffer overflow, and SQL-class attacks
are architecturally impossible rather than merely mitigated. We define
the \emph{Tempus} language with a formal BNF grammar, small-step
operational semantics, and a compile-time cell registry that is provably
immutable at runtime. COMPILE and EXECUTE phases are shown to be
mutually exclusive states of the trajectory state machine. Security
theorems with formal proofs cover payload injection impossibility, replay
failure under monotone cycle counts, and bounded attack surface. We
discuss hardware realizations using crystal, MEMS, and atomic oscillators,
and present applications in nuclear process control, autonomous drone
systems, and industrial IoT networks.
\end{abstract}
\vspace{1em}

\end{@twocolumnfalse}
}]

%=======================================================================
\section{Introduction}
\label{sec:intro}
%=======================================================================

\subsection{Motivation}

Modern distributed systems rely on rich message protocols: TCP/IP
packets carry headers, payloads, checksums, and framing bytes that are
parsed at every layer of the network stack. This richness is both a
strength and an inexhaustible attack surface. Buffer overflows exploit
the boundary between parsing code and memory allocators
\cite{saltzer1975protection}. SQL injection exploits the conflation of
code and data in database query engines. Command injection exploits
shell interpreters that treat user-supplied strings as executable text.
Despite fifty years of mitigations --- bounds checking, memory-safe
languages, sandboxing, formal verification --- payload-based attacks
remain a leading cause of system compromise. The mitigations address
symptoms while preserving the root cause: the existence of a parser.

Real-time and safety-critical distributed systems face a compounded
version of this problem. A nuclear control system, an aircraft fly-by-wire
network, or an autonomous drone swarm must not only be
\emph{functionally correct} but \emph{structurally immune} to adversarial
manipulation. Formal verification can establish correctness of a
protocol implementation, but cannot prevent injection of malformed
messages that exercise unverified corner cases. The fundamental
difficulty is that content-addressed protocols endow every message with
expressive power that may be abused.

\subsection{The Temporal Programming Hypothesis}

We propose a different foundation. Suppose every inter-node signal
is a bare pulse: a single edge with no amplitude encoding, no length,
no header. The only observable property is \emph{when} it arrives
relative to the local oscillator. The datum is:
\begin{equation}
  \DP(k) \;=\; \Tref(k) - \trec(k)
  \label{eq:dp-def}
\end{equation}
where $\Tref(k)$ is the time predicted by the local reference oscillator
for the $k$-th pulse and $\trec(k)$ is its observed arrival time.
Programs are written as \emph{timing cells}: measurable intervals in
$\DP$-space. Execution is the assignment of an observed $\DP(k)$ to its
enclosing cell, followed by dispatch of the cell's pre-compiled action.

Nothing in this scheme requires a parser. There is no payload to
buffer, no string to interpret, no code to inject. The attack surface
is not reduced --- it is \emph{absent by construction}.

\subsection{Contributions}

This paper makes the following original contributions:

\begin{enumerate}[leftmargin=1.5em,topsep=2pt,itemsep=1pt]
  \item \textbf{Temporal programming model.} We formalize the
    oscillator model, the $\DP$ primitive, timing cells as measurable
    sets, and labeled trajectory composition
    (Section~\ref{sec:model}).

  \item \textbf{Tempus language.} We define a minimal temporal
    programming language with BNF grammar, small-step operational
    semantics, and worked examples
    (Section~\ref{sec:tempus}).

  \item \textbf{Composition Inflation Theorem.} We prove
    $T(n,d) = d\,(1+d)^{n-1}$ from first principles via
    the stars-and-bars bijection and the binomial theorem
    (Section~\ref{sec:inflation}).

  \item \textbf{Formal security analysis.} We prove payload injection
    impossibility, replay failure, and bounded attack surface under a
    formal adversary model (Section~\ref{sec:security}).

  \item \textbf{Implementation architecture.} We give pseudocode for
    the cell registry, trajectory engine, and scheduler
    (Section~\ref{sec:impl}).
\end{enumerate}

\subsection{Organization}

Section~\ref{sec:background} reviews related paradigms.
Section~\ref{sec:model} develops the temporal programming model.
Section~\ref{sec:tempus} defines the Tempus language.
Section~\ref{sec:inflation} proves the Composition Inflation Theorem.
Section~\ref{sec:security} gives the security analysis.
Section~\ref{sec:impl} describes the implementation architecture.
Section~\ref{sec:apps} discusses applications.
Section~\ref{sec:eval} evaluates complexity and overhead.
Section~\ref{sec:discussion} discusses limitations and future work.
Section~\ref{sec:conclusion} concludes.

%=======================================================================
\section{Background}
\label{sec:background}
%=======================================================================

\subsection{Event-Triggered and Time-Triggered Systems}

Two canonical paradigms govern real-time distributed computation.
In \emph{event-triggered} systems, a node transmits a message when
a monitored condition changes. The message carries the new value.
This model is reactive and efficient under light load but susceptible
to message storms under overload \cite{kopetz2003tta}.

In \emph{time-triggered} systems, pioneered by Kopetz in the
Time-Triggered Architecture (TTA), communication is governed by a
static schedule: nodes transmit in pre-allocated time slots
\cite{kopetz2003tta,kopetz2011rts}. TTA eliminates message storms
and provides deterministic latency bounds, but still transmits
content-carrying frames that must be parsed.

Temporal programming occupies a third position: the schedule is
implicit in the oscillator frequency, and the \emph{only} information
is whether each pulse arrives early, on time, or late. No frame is
transmitted; no parser is required.

\subsection{Synchronous Programming Languages}

The synchronous hypothesis \cite{berry1992esterel} states that a
program's reaction to an input is logically instantaneous relative
to the environment's clock. \emph{Esterel} \cite{berry1992esterel}
and \emph{Lustre} \cite{halbwachs1991lustre} build on this hypothesis
to achieve deterministic concurrency. Both languages operate on streams
of typed values and require a parser for message content.

Temporal programming is synchronous in a deeper sense: the reference
clock is not an abstraction but a physical oscillator, and the
``value'' of every event is its timing relative to that oscillator.
There is no value stream to parse because timing deviation is the
value.

\subsection{Clock Synchronization}

Lamport's seminal work \cite{lamport1978clocks} established that
without a global clock, event ordering in a distributed system must
be inferred from happened-before relations. The Network Time Protocol
\cite{mills1991ntp} and the IEEE~1588 Precision Time Protocol
\cite{ieee1588} provide sub-microsecond synchronization over IP
networks. Mattern \cite{mattern1989virtualtime} generalized Lamport's
approach to virtual time. In temporal programming, synchronization
quality directly bounds cell resolution: a tighter synchronized clock
allows narrower cells (Section~\ref{sec:model}).

\subsection{Temporal Logic and Timed Automata}

Pnueli's temporal logic \cite{pnueli1977temporal} provides a
specification language for time-dependent properties of programs.
Alur and Dill's timed automata \cite{alur1994timed} extend finite
automata with real-valued clock variables and guards. Henzinger's
hybrid automata \cite{henzinger1996hybrid} further extend to
continuous dynamics. These frameworks specify \emph{when} a system
should be in a given state; temporal programming specifies
\emph{what} timing deviations trigger which actions. The relationship
is that a Tempus program can be compiled to a timed automaton for
model-checking purposes.

\subsection{Sensor Networks and PTIDES}

TinyOS \cite{levis2005tinyos} demonstrated that resource-constrained
sensor nodes can execute sophisticated data-collection programs.
Hill et al.\ \cite{hill2000sensor} showed that networked sensor
architectures must co-design hardware and software. PTIDES
\cite{zhao2007ptides} provides a programming model for
time-synchronized distributed real-time systems based on the
Ptolemy~II framework. Temporal programming differs: PTIDES
still carries typed data values; temporal programming eliminates
the value entirely.

%=======================================================================
\section{The Temporal Programming Model}
\label{sec:model}
%=======================================================================

\subsection{Oscillator Model}

\begin{definition}[Reference Oscillator]
\label{def:oscillator}
A \emph{reference oscillator} is a tuple $\mathcal{O} =
(\fref, \epsilon, M_0)$ where:
\begin{itemize}[leftmargin=1.5em,topsep=2pt,itemsep=1pt]
  \item $\fref \in \mathbb{R}_{>0}$ is the nominal frequency in Hz,
  \item $\epsilon \in [0,1)$ is the maximum fractional frequency
    deviation (stability bound),
  \item $M_0 \in \mathbb{N}$ is the initial cycle count at deployment.
\end{itemize}
The oscillator generates a strictly increasing sequence of
\emph{reference times}
$\{\Tref(k)\}_{k \geq M_0}$ defined by:
\begin{equation}
  \Tref(k) = \frac{k}{\fref}
  \quad \forall k \in \mathbb{N},\; k \geq M_0.
  \label{eq:tref}
\end{equation}
The \emph{cycle count} $\Mcyc(t)$ at physical time $t$ is the
(integer) number of complete oscillator periods elapsed since $M_0$:
\begin{equation}
  \Mcyc(t) = M_0 + \lfloor (t - \Tref(M_0))\,\fref \rfloor.
  \label{eq:cyclecnt}
\end{equation}
\end{definition}

\begin{lemma}[Monotonicity of Cycle Count]
\label{lem:mono}
For all physical times $t_1 < t_2$, $\Mcyc(t_1) \leq \Mcyc(t_2)$.
\end{lemma}
\begin{proof}
By definition~\eqref{eq:cyclecnt},
$\Mcyc(t) = M_0 + \lfloor (t - \Tref(M_0))\,\fref \rfloor$.
Since $\fref > 0$ and the floor function is non-decreasing,
$t_1 < t_2 \Rightarrow (t_1 - \Tref(M_0))\,\fref <
(t_2 - \Tref(M_0))\,\fref$, so
$\lfloor (t_1 - \Tref(M_0))\,\fref \rfloor \leq
\lfloor (t_2 - \Tref(M_0))\,\fref \rfloor$.
Hence $\Mcyc(t_1) \leq \Mcyc(t_2)$. \qed
\end{proof}

\subsection{The Timing Deviation Primitive}

\begin{definition}[Timing Deviation]
\label{def:dp}
Let $\mathcal{O}$ be a reference oscillator. For the $k$-th received
pulse with actual reception time $\trec(k)$, the \emph{timing
deviation} is:
\begin{equation}
  \DP(k) \;=\; \Tref(k) - \trec(k) \;\in\; \mathbb{R}.
  \label{eq:dp}
\end{equation}
$\DP(k) > 0$ indicates an early arrival (pulse precedes the reference
tick); $\DP(k) < 0$ indicates a late arrival; $\DP(k) = 0$ indicates
perfect synchrony.
\end{definition}

The $\DP$-space is $\mathbb{R}$. In practice, given maximum clock skew
$\epsilon$ and propagation jitter bound $\delta$, all physically
realizable deviations lie in the bounded interval:
\begin{equation}
  \calD \;=\; \left[
    -\frac{\epsilon + \delta}{\fref},\;
    \frac{\epsilon + \delta}{\fref}
  \right] \subset \mathbb{R}.
  \label{eq:dpspace}
\end{equation}
We henceforth work within $\calD$.

\subsection{Timing Cells}

\begin{definition}[Timing Cell]
\label{def:cell}
A \emph{timing cell} over $\calD$ is a Borel-measurable set
$\calC \subseteq \calD$ with positive Lebesgue measure:
$\Leb(\calC) > 0$. A \emph{cell partition} of $\calD$ is a finite
collection $\calP = \{\calC_1, \ldots, \calC_m\}$ such that:
\begin{enumerate}[leftmargin=1.5em,topsep=2pt,itemsep=1pt]
  \item[(i)]   $\bigcup_{i=1}^{m} \calC_i = \calD$
    (coverage),
  \item[(ii)]  $\calC_i \cap \calC_j = \emptyset$ for $i \neq j$
    (disjointness, up to measure-zero boundaries),
  \item[(iii)] $\Leb(\calC_i) > 0$ for all $i$
    (non-degeneracy).
\end{enumerate}
\end{definition}

The positive-measure requirement is essential: it ensures that
$\calC_i$ occupies a physical interval of time and is not a point
predicate. Point predicates have Lebesgue measure zero and cannot
be robustly detected in the presence of oscillator jitter.

\begin{figure*}[!htbp]
    \centering
    \includegraphics[width=\textwidth]{figures/panel_02.png}
    \caption{\textbf{Timing Cell and S-Functional Floor:
    operational truth is a cell $C \subseteq \Delta P$-space with
    positive measure, and the S-functional satisfies
    $S(R,x;\,C) \geq \beta > 0$ for all $x$.}
    (\textbf{A}) Histogram of 4000 simulated $\Delta P$ values
    (blue bars, density-normalised) with the timing cell $C$
    bounded by red dashed lines.
    The finite-width interval captures a substantial fraction of the
    distribution; no point predicate ($\Delta P = c$) could do so
    with positive measure.
    (\textbf{B}) S-functional $S(R,x;\,C)$ versus position $x$ for
    a one-dimensional cell $C = [0.3,\,0.7]$ (blue curve).
    Outside $C$ the function is the V-shaped distance to the
    nearest cell boundary; inside $C$ it is flat at the floor $\beta$
    (red dashed line, green shading), confirming that
    $\inf_{x} S = \beta > 0$ strictly.
    (\textbf{C}) Three-dimensional S-functional surface over a
    two-dimensional $\Delta P$ domain, where the cell is the square
    $[0.3,0.7]^2$.
    The surface is a four-sided bowl whose flat base lies at $\beta$;
    the gradient outside the cell encodes the distance to the nearest
    cell face.
    (\textbf{D}) Effective measure $\tau(C) + \beta$ versus cell width
    for four values of the floor parameter $\beta$
    ($0.02$, $0.05$, $0.10$, $0.20$).
    All curves are strictly positive for all cell widths, demonstrating
    that the floor $\beta$ provides a hard lower bound independent of
    cell geometry.}
    \label{fig:cell_floor}
\end{figure*}

\begin{definition}[Action-Equivalence]
\label{def:actequiv}
Two timing deviations $x, y \in \calD$ are \emph{action-equivalent
under partition $\calP$}, written $x \sim_{\calP} y$, iff there
exists $\calC_i \in \calP$ with $x \in \calC_i$ and $y \in \calC_i$.
\end{definition}

\begin{proposition}[Equivalence Relation]
$\sim_{\calP}$ is an equivalence relation on $\calD$.
\end{proposition}
\begin{proof}
Reflexivity: $x \sim_{\calP} x$ because coverage ensures $x$ lies
in some $\calC_i$. Symmetry: if $x, y \in \calC_i$ then
$y, x \in \calC_i$. Transitivity: if $x, y \in \calC_i$ and
$y, z \in \calC_j$, then since the partition is disjoint and $y \in
\calC_i \cap \calC_j$, we have $i = j$, so $x, z \in \calC_i$. \qed
\end{proof}

\begin{theorem}[Cell Action-Equivalence]
\label{thm:cellequiv}
Let $\mathcal{A} : \calP \to \mathcal{F}$ be a cell-action map
assigning a function $f_i \in \mathcal{F}$ to each cell $\calC_i$.
Then $x \sim_{\calP} y \Rightarrow \mathcal{A}(\cell(x)) =
\mathcal{A}(\cell(y))$, where $\cell(x)$ is the unique cell
containing $x$.
\end{theorem}
\begin{proof}
If $x \sim_{\calP} y$ then by Definition~\ref{def:actequiv} there
exists $\calC_i$ with $x, y \in \calC_i$, so
$\cell(x) = \cell(y) = \calC_i$, hence
$\mathcal{A}(\cell(x)) = \mathcal{A}(\calC_i) = \mathcal{A}(\cell(y))$.
\qed
\end{proof}

Theorem~\ref{thm:cellequiv} formalizes the core semantic principle:
\emph{what matters is which cell a deviation falls in, not its
precise value}.

\subsection{Labeled Channels and Trajectories}

Physical systems emit signals from multiple sources: temperature
sensors, pressure sensors, actuator feedback, etc. Each source
is modeled as a distinct channel.

\begin{definition}[Channel]
\label{def:channel}
A \emph{channel} is a pair $\chi = (s, \mathcal{O}_s)$ where $s$ is
a source identifier and $\mathcal{O}_s$ is the reference oscillator
for that source. The set of active channels is
$\{\chi_1, \ldots, \chi_d\}$ with $d \geq 1$.
\end{definition}

\begin{definition}[Labeled Timing Event]
\label{def:event}
A \emph{labeled timing event} is a pair $e = (\chi_i, \DP)$ where
$\chi_i$ is the channel and $\DP \in \calD_i$ is the timing
deviation on that channel. The event time is implicitly
$\trec = T_{\mathrm{ref},i} - \DP$.
\end{definition}

\begin{definition}[Trajectory]
\label{def:traj}
A \emph{trajectory of length $n$} is an ordered sequence of labeled
timing events:
\begin{equation}
  \tau = (e_1, e_2, \ldots, e_n) =
  ((\chi_{i_1}, \DP_1), \ldots, (\chi_{i_n}, \DP_n)).
  \label{eq:traj}
\end{equation}
The \emph{channel sequence} of $\tau$ is
$\lambda(\tau) = (\chi_{i_1}, \ldots, \chi_{i_n}) \in
\{\chi_1, \ldots, \chi_d\}^n$.
Two trajectories are \emph{cell-equivalent} if each corresponding
pair of deviations is action-equivalent under the respective channel
partition.
\end{definition}

\subsection{Multi-Domain Physical Transduction}

A central claim of temporal programming is that arbitrary physical
quantities reduce to timing deviations.

\begin{definition}[Physical Transducer]
\label{def:transducer}
A \emph{physical transducer for quantity $Q$} is a device that
converts a measurement $q$ of $Q$ to a timing deviation
$\DP_Q(q) \in \calD_Q \subset \mathbb{R}$.
\end{definition}

\begin{example}[Temperature via Crystal Oscillator]
\label{ex:temp}
A temperature-compensated crystal oscillator (TCXO) has frequency
$f(T) = f_0(1 + \alpha(T - T_0) + \beta(T - T_0)^2)$ where $T$ is
temperature and $\alpha, \beta$ are crystal-specific coefficients.
At temperature $T$, the $k$-th pulse arrives at
$\trec(k; T) = k / f(T)$, producing
\begin{equation}
  \DP_{\mathrm{temp}}(k; T)
  = \frac{k}{\fref} - \frac{k}{f(T)}
  = k \left(\frac{1}{\fref} - \frac{1}{f(T)}\right),
  \label{eq:dp-temp}
\end{equation}
which is a monotone function of $T$ for fixed $k$, $\fref$. Each
temperature range $[T_a, T_b]$ maps to a cell in $\calD_{\mathrm{temp}}$.
\end{example}

\begin{example}[Pressure via MEMS Resonator]
\label{ex:press}
A MEMS resonant pressure sensor shifts its resonant period as
$\tau_{\mathrm{MEMS}}(P) = \tau_0 / (1 + \gamma P)$,
where $P$ is applied pressure and $\gamma$ is the pressure
coefficient. Measuring the elapsed time over $k$ MEMS cycles
gives $\DP_{\mathrm{press}} = k(\tau_0 - \tau_{\mathrm{MEMS}}(P))$,
again a monotone map to $\calD_{\mathrm{press}}$.
\end{example}

\begin{proposition}[Multi-Domain Composition]
\label{prop:multidom}
Let $Q_1, \ldots, Q_d$ be physical quantities with respective
transducers $\DP_{Q_1}, \ldots, \DP_{Q_d}$ and $\DP$-spaces
$\calD_1, \ldots, \calD_d$. The \emph{composed trajectory space} is
$\calD_1 \times \cdots \times \calD_d$ with channel label
$\chi_i$ indexing the $i$-th coordinate. A $d$-channel trajectory
$\tau$ of length $n$ is a path in this product space.
\end{proposition}
\begin{proof}
Each event $e = (\chi_i, \DP_i)$ contributes an element of $\calD_i$.
The channel label $\chi_i$ uniquely identifies the coordinate.
The Cartesian product $\prod_{i=1}^{d} \calD_i$ is well-defined as a
measurable space under the product $\sigma$-algebra, and trajectories
are sequences of elements drawn from this product space. \qed
\end{proof}

%=======================================================================
\section{The Tempus Language}
\label{sec:tempus}
%=======================================================================

\subsection{Abstract Syntax}

We define Tempus via a BNF grammar. Meta-variables: $x$ ranges over
identifiers; $r \in \mathbb{R}$ over real literals; $n \in \mathbb{N}$
over natural number literals.

\begin{figure}[ht]
\small
\begin{tabbing}
\hspace{1em}\=\hspace{0.5em}\=\kill
$\langle\mathit{prog}\rangle$  \> $::=$ \> $\langle\mathit{decl}\rangle^*$  \\[2pt]
$\langle\mathit{decl}\rangle$  \> $::=$ \>
  $\langle\mathit{cell\text{-}decl}\rangle$
  $\mid$ $\langle\mathit{sync\text{-}decl}\rangle$
  $\mid$ $\langle\mathit{compose\text{-}decl}\rangle$
  $\mid$ $\langle\mathit{when\text{-}decl}\rangle$
  $\mid$ $\langle\mathit{ensemble\text{-}decl}\rangle$\\[2pt]
$\langle\mathit{cell\text{-}decl}\rangle$ \> $::=$ \>
  \textbf{cell} $x$
  \textbf{bounds} $(r, r)$
  \textbf{action} $\langle\mathit{expr}\rangle$ \\[2pt]
$\langle\mathit{sync\text{-}decl}\rangle$ \> $::=$ \>
  \textbf{sync} $x$
  \textbf{at} $r$ \textbf{freq} \\[2pt]
$\langle\mathit{compose\text{-}decl}\rangle$ \> $::=$ \>
  \textbf{compose} $\mathbf{d}{=}n$
  \textbf{channels} $x, \ldots, x$
  \textbf{into} $x$ \\[2pt]
$\langle\mathit{when\text{-}decl}\rangle$ \> $::=$ \>
  \textbf{when} $x$ \textbf{do} $\langle\mathit{stmt}\rangle$ \\[2pt]
$\langle\mathit{ensemble\text{-}decl}\rangle$ \> $::=$ \>
  \textbf{ensemble} $x^+$
  \textbf{sync\_threshold} $r$ \\[2pt]
$\langle\mathit{expr}\rangle$ \> $::=$ \>
  $r$ $\mid$ $x$
  $\mid$ $\langle\mathit{expr}\rangle$ $\mathit{binop}$
  $\langle\mathit{expr}\rangle$
  $\mid$ \textbf{let} $x = \langle\mathit{expr}\rangle$
  \textbf{in} $\langle\mathit{expr}\rangle$
  $\mid$ \textbf{if} $\langle\mathit{bexpr}\rangle$
  \textbf{then} $\langle\mathit{expr}\rangle$
  \textbf{else} $\langle\mathit{expr}\rangle$ \\[2pt]
$\langle\mathit{stmt}\rangle$ \> $::=$ \>
  \textbf{emit} $x$
  $\mid$ \textbf{wait} $r$
  $\mid$ \textbf{fire} $\langle\mathit{expr}\rangle$
  $\mid$ $\langle\mathit{stmt}\rangle$ \textbf{;} $\langle\mathit{stmt}\rangle$
  $\mid$ \textbf{begin} $\langle\mathit{stmt}\rangle^+$ \textbf{end} \\[2pt]
$\langle\mathit{bexpr}\rangle$ \> $::=$ \>
  \textbf{true} $\mid$ \textbf{false}
  $\mid$ $\langle\mathit{expr}\rangle$ $\mathit{cmp}$
  $\langle\mathit{expr}\rangle$
  $\mid$ $\langle\mathit{bexpr}\rangle$ \textbf{and}
  $\langle\mathit{bexpr}\rangle$
  $\mid$ \textbf{not} $\langle\mathit{bexpr}\rangle$
\end{tabbing}
\caption{Abstract syntax of Tempus.}
\label{fig:bnf}
\end{figure}

\subsection{Static Semantics and Typing}

A Tempus program is \emph{well-formed} if:
(i) every \textbf{when} clause references a declared cell name;
(ii) every \textbf{compose} clause references declared sync sources;
(iii) every \textbf{ensemble} clause references declared agents;
(iv) for each channel, the cell bounds tile $\calD$ without gap
  (coverage check).
The compiler rejects any program that fails these checks before
emitting any runtime artifact.

\subsection{Operational Semantics}

We define a small-step reduction relation $\to$ on configurations.
A \emph{configuration} is a tuple $\langle \sigma, \Gamma, \tau,
\phi \rangle$ where:
\begin{itemize}[leftmargin=1.5em,topsep=2pt,itemsep=1pt]
  \item $\sigma \in \{\mathsf{COMPILE},\mathsf{EXECUTE},\mathsf{IDLE}\}$
    is the phase,
  \item $\Gamma$ is the cell registry (compile-time constant),
  \item $\tau$ is the current (open or closed) trajectory,
  \item $\phi$ is the pending statement queue.
\end{itemize}

\textbf{Rules:}

\begin{align}
  &\textsc{(Receive)} \quad
  \frac{e = (\chi, \DP) \quad \sigma = \mathsf{COMPILE}}
       {\langle \mathsf{COMPILE},\Gamma,\tau,\phi\rangle
        \;\to\;
        \langle \mathsf{COMPILE},\Gamma,\tau \cdot e,\phi\rangle}
  \label{rule:recv} \\[6pt]
  &\textsc{(Assign)} \quad
  \frac{|\tau| = n \quad \calC = \traj(\tau, \Gamma) \neq \bot}
       {\langle \mathsf{COMPILE},\Gamma,\tau,\phi\rangle
        \;\to\;
        \langle \mathsf{EXECUTE},\Gamma,\tau,\Gamma(\calC)\rangle}
  \label{rule:assign} \\[6pt]
  &\textsc{(Fire)} \quad
  \frac{\sigma = \mathsf{EXECUTE} \quad \phi = s;\phi'}
       {\langle \mathsf{EXECUTE},\Gamma,\tau,s;\phi'\rangle
        \;\to\;
        \langle \mathsf{EXECUTE},\Gamma,\tau,\phi'\rangle \;\|\;
        \llb s \rrb}
  \label{rule:fire} \\[6pt]
  &\textsc{(Reset)} \quad
  \frac{\sigma = \mathsf{EXECUTE} \quad \phi = \varepsilon}
       {\langle \mathsf{EXECUTE},\Gamma,\tau,\varepsilon\rangle
        \;\to\;
        \langle \mathsf{IDLE},\Gamma,(),\varepsilon\rangle}
  \label{rule:reset} \\[6pt]
  &\textsc{(Idle)} \quad
  \frac{\sigma = \mathsf{IDLE}}
       {\langle \mathsf{IDLE},\Gamma,(),\varepsilon\rangle
        \;\to\;
        \langle \mathsf{COMPILE},\Gamma,(),\varepsilon\rangle}
  \label{rule:idle}
\end{align}

The $\traj(\tau, \Gamma)$ function maps a complete trajectory $\tau$
of length $n$ to a cell by looking up each event's $\DP$ in the
corresponding channel partition in $\Gamma$ and returning the
composite cell identifier. $\llb s \rrb$ denotes the denotational
semantics of statement $s$ (side effects in the environment).

\subsection{COMPILE/EXECUTE Phase Separation}

\begin{theorem}[Phase Mutual Exclusion]
\label{thm:phase}
In any execution trace, the phases $\mathsf{COMPILE}$ and
$\mathsf{EXECUTE}$ never hold simultaneously in the same
configuration.
\end{theorem}
\begin{proof}
Each configuration carries exactly one phase tag
$\sigma \in \{\mathsf{COMPILE}, \mathsf{EXECUTE}, \mathsf{IDLE}\}$.
No reduction rule produces a configuration with two simultaneous
phase tags: rules \eqref{rule:recv} and \eqref{rule:assign}
emit a single phase value in their conclusion, as do rules
\eqref{rule:fire}, \eqref{rule:reset}, and \eqref{rule:idle}.
The tag is a function, not a set. Hence at any configuration in any
execution trace, $\sigma$ takes exactly one value, and the events
$\sigma = \mathsf{COMPILE}$ and $\sigma = \mathsf{EXECUTE}$ are
mutually exclusive. \qed
\end{proof}

\begin{corollary}[Semantic Justification of Phase Separation]
\label{cor:phaselogic}
$\mathsf{COMPILE}$ requires an open trajectory (undetermined cell
assignment); $\mathsf{EXECUTE}$ requires a closed trajectory
(determined cell assignment). These conditions are contradictory:
\begin{enumerate}[leftmargin=1.5em,topsep=2pt,itemsep=1pt]
  \item[$\bullet$] An open trajectory has $|\tau| < n$, so
    $\traj(\tau, \Gamma) = \bot$: the cell is unknown.
  \item[$\bullet$] Rule \eqref{rule:assign} transitions to
    $\mathsf{EXECUTE}$ only when $|\tau| = n$ and
    $\traj(\tau, \Gamma) \neq \bot$: the cell is known.
\end{enumerate}
These preconditions are logically incompatible, providing a second,
semantic proof of mutual exclusion. \qed
\end{corollary}

\begin{figure*}[!htbp]
    \centering
    \includegraphics[width=\textwidth]{figures/panel_04.png}
    \caption{\textbf{COMPILE/EXECUTE Phase Separation Forced by the
    Receiver Uncertainty Principle
    $\sigma_K \cdot \sigma_Y \geq \hbar_R = \beta\tau$:
    trajectory accumulation and action firing are mutually exclusive
    phases.}
    (\textbf{A}) Antiphase oscillation of $\sigma_K$ (knowledge
    uncertainty, blue), $\sigma_Y$ (action uncertainty, red), and
    their product $\sigma_K \cdot \sigma_Y$ (green dotted) over four
    complete cycles.
    The product never drops below the floor $\hbar_R$ (grey dashed),
    and $\sigma_K$ is maximised precisely when $\sigma_Y$ is minimised,
    and vice versa.
    (\textbf{B}) Timeline of alternating COMPILE (blue) and EXECUTE
    (red) blocks over ten cycles.
    Block durations are drawn from realistic distributions; the two
    phases tile the time axis without overlap, illustrating the
    hard mutual exclusion enforced by the uncertainty principle.
    (\textbf{C}) Three-dimensional action-rate surface over
    $(\sigma_K,\,\sigma_Y) \in [0,1]^2$.
    The surface is nonzero only in the region $\sigma_K < 0.2$
    (deep EXECUTE phase) and grows with $\sigma_Y$, reflecting that
    actions fire at high confidence only after the trajectory has
    been fully resolved.
    (\textbf{D}) Action counts per block for 120 alternating blocks.
    COMPILE blocks (blue) fire zero actions by construction; EXECUTE
    blocks (red) fire a Poisson-distributed count $\sim\mathrm{Poi}(3.5)$,
    confirming that no action escapes the phase boundary.}
    \label{fig:phase_separation}
\end{figure*}

\subsection{Tempus Code Examples}

\begin{example}[Nuclear Coolant Temperature Monitor]
\label{ex:nuclear}
The following Tempus program monitors reactor coolant temperature
via a TCXO-based transducer and dispatches alerts.
\begin{lstlisting}[language=Tempus, caption={Nuclear coolant monitor in Tempus.}, label=lst:nuclear]
-- Declare reference oscillator
sync coolant_sensor at 10.0e6 freq

-- Timing cells for coolant temperature bands
cell NOMINAL  bounds (-1.0e-7,  1.0e-7) action 0
cell WARM     bounds ( 1.0e-7,  5.0e-7) action 1
cell HOT      bounds ( 5.0e-7,  2.0e-6) action 2
cell CRITICAL bounds ( 2.0e-6, 10.0e-6) action 3

-- Trajectory: 4 events, single channel
compose d=1 channels coolant_sensor into coolant_traj

-- Dispatch rules
when NOMINAL  do emit status_ok
when WARM     do emit status_warn
when HOT      do begin
                   emit status_hot;
                   fire reduce_power(0.8)
                 end
when CRITICAL do begin
                   emit scram_alert;
                   fire emergency_shutdown
                 end
\end{lstlisting}
\end{example}

\begin{example}[Drone Altitude-Velocity Compose]
\label{ex:drone}
A two-channel composition tracks both altitude (barometric) and
vertical velocity (IMU) simultaneously.
\begin{lstlisting}[language=Tempus, caption={Drone dual-channel trajectory.}, label=lst:drone]
sync baro_sensor   at 1.0e6 freq
sync imu_sensor    at 4.0e6 freq

-- Altitude cells
cell ALT_LOW    bounds (-5.0e-6, -1.0e-6) action -1
cell ALT_HOLD   bounds (-1.0e-6,  1.0e-6) action  0
cell ALT_HIGH   bounds ( 1.0e-6,  5.0e-6) action  1

-- Velocity cells
cell VEL_DESCEND bounds (-3.0e-6, -0.5e-6) action -1
cell VEL_STABLE  bounds (-0.5e-6,  0.5e-6) action  0
cell VEL_CLIMB   bounds ( 0.5e-6,  3.0e-6) action  1

-- 6-event, 2-channel composition
compose d=2 channels baro_sensor, imu_sensor
        into altitude_velocity_space

-- Ensemble of 4 drones with 0.1 ms sync threshold
ensemble drone_1, drone_2, drone_3, drone_4
         sync_threshold 1.0e-4

when ALT_LOW  do fire thrust_up(0.3)
when ALT_HOLD do fire hold_throttle
when ALT_HIGH do fire thrust_down(0.3)
\end{lstlisting}
\end{example}

%=======================================================================
\section{Trajectory Distinguishability: The Composition Inflation Theorem}
\label{sec:inflation}
%=======================================================================

\subsection{Counting Labeled Compositions}

We now prove the central combinatorial result. The question is:
given $n$ timing events distributed across $d$ active channels,
how many categorically distinct labeled trajectories can the
system distinguish?

\begin{definition}[Labeled Composition]
\label{def:labeledcomp}
A \emph{composition of $n$} is an ordered sequence of positive integers
$(c_1, c_2, \ldots, c_k)$ with $c_1 + c_2 + \cdots + c_k = n$ and
$k \geq 1$. A \emph{$d$-labeled composition of $n$} assigns to each
part $c_j$ a label $\ell_j \in \{1, \ldots, d\}$.
\end{definition}

\begin{theorem}[Composition Inflation]
\label{thm:inflation}
For $n \geq 1$ timing events with $d \geq 1$ active channels, the
number of categorically distinguishable $d$-labeled compositions of
$n$ is:
\begin{equation}
  T(n, d) \;=\; d\,(1 + d)^{n-1}.
  \label{eq:inflation}
\end{equation}
\end{theorem}
\begin{proof}
The proof proceeds in three steps.

\medskip
\noindent\textbf{Step 1: Counting unlabeled compositions.}
A composition of $n$ into $k$ positive parts corresponds bijectively
to a choice of $k-1$ dividers placed among the $n-1$ gaps between $n$
indistinguishable dots (the stars-and-bars bijection
\cite{andrews1976partitions}). The number of compositions of $n$ into
exactly $k$ parts is $\binom{n-1}{k-1}$. Summing over all $k$:
\begin{equation}
  \sum_{k=1}^{n} \binom{n-1}{k-1} \;=\; \sum_{j=0}^{n-1} \binom{n-1}{j}
  \;=\; 2^{n-1}.
  \label{eq:total-comps}
\end{equation}
Thus there are $2^{n-1}$ compositions of $n$ in total, a classical
result \cite{graham1994concrete}.

\medskip
\noindent\textbf{Step 2: Counting $d$-labeled compositions.}
For a composition with exactly $k$ parts, each part independently
receives one of $d$ channel labels. There are $d^k$ such label
assignments. The total count of $d$-labeled compositions of $n$ is
therefore:
\begin{equation}
  T(n,d) \;=\; \sum_{k=1}^{n} \binom{n-1}{k-1} \cdot d^k.
  \label{eq:Tnd-sum}
\end{equation}

\medskip
\noindent\textbf{Step 3: Closed form via the binomial theorem.}
Factor out $d$:
\begin{equation}
  T(n,d) \;=\; d \sum_{k=1}^{n} \binom{n-1}{k-1} d^{k-1}
  \;=\; d \sum_{j=0}^{n-1} \binom{n-1}{j} d^{j}.
  \label{eq:factor}
\end{equation}
The inner sum is the binomial expansion of $(1+d)^{n-1}$ by the
binomial theorem \cite{graham1994concrete}:
\begin{equation}
  \sum_{j=0}^{n-1} \binom{n-1}{j} d^{j}
  \;=\; (1 + d)^{n-1}.
  \label{eq:binom-app}
\end{equation}
Substituting:
\begin{equation}
  T(n,d) = d\,(1+d)^{n-1}.
  \label{eq:Tnd-closed}
\end{equation}
\qed
\end{proof}

\begin{corollary}[Growth Rate]
$T(n,d) = \Theta(d^n)$ as $n \to \infty$ for fixed $d \geq 2$, and
$T(n,d) = \Theta(n\,d^{n-1})$ differentiating in $d$ for fixed $n$.
The state space is exponential in the trajectory length.
\end{corollary}

\begin{corollary}[Special Cases]
\label{cor:special}
\begin{enumerate}[leftmargin=1.5em,topsep=2pt,itemsep=1pt]
  \item $T(n,1) = 1 \cdot 2^{n-1} = 2^{n-1}$: single-channel system,
    $2^{n-1}$ distinct trajectories.
  \item $T(1, d) = d \cdot (1+d)^0 = d$: single event, $d$ channels,
    trivially $d$ trajectories (one per channel).
  \item $T(n, d) > T(n, d-1)$ for all $n \geq 1$, $d \geq 2$:
    adding a channel strictly inflates the trajectory count.
\end{enumerate}
\end{corollary}

\begin{table}[ht]
\centering
\caption{Values of $T(n,d) = d\,(1+d)^{n-1}$}
\label{tab:tnd}
\small
\begin{tabular}{@{}lrrrrr@{}}
\toprule
$n$ \textbackslash $d$ & 1 & 2 & 3 & 4 & 8 \\
\midrule
1  &   1 &    2 &    3 &    4 &    8 \\
2  &   2 &    6 &   12 &   20 &   72 \\
3  &   4 &   18 &   48 &  100 &  648 \\
4  &   8 &   54 &  192 &  500 & 5832 \\
5  &  16 &  162 &  768 & 2500 & 52\,488 \\
6  &  32 &  486 & 3072 &12\,500 & 472\,392 \\
\bottomrule
\end{tabular}
\end{table}

\subsection{Monotone Partition Count and Replay Immunity}

\begin{theorem}[Monotone Partition Count]
\label{thm:mono-partition}
Let $M_0 \leq M_1$ be two cycle counts with $M_0 \neq M_1$.
Let a signal be emitted at cycle $M_0$ with timing deviation
$\DP_0 = \Tref(M_0) - t_{\mathrm{rec},0}$. If the same physical signal
is replayed at cycle $M_1 > M_0$, the resulting timing deviation
is $\DP_1 = \Tref(M_1) - t_{\mathrm{rec},0} \neq \DP_0$.
\end{theorem}
\begin{proof}
By Definition~\ref{def:oscillator},
$\Tref(M) = M / \fref$ for integer $M$.
\begin{align}
  \DP_0 &= M_0/\fref - t_{\mathrm{rec},0}, \label{eq:dp0} \\
  \DP_1 &= M_1/\fref - t_{\mathrm{rec},0}. \label{eq:dp1}
\end{align}
Subtracting: $\DP_1 - \DP_0 = (M_1 - M_0)/\fref$.
Since $M_1 > M_0$ and $\fref > 0$, we have $\DP_1 - \DP_0 > 0$,
so $\DP_1 \neq \DP_0$. \qed
\end{proof}

The implication is that a recorded timing deviation from an earlier
epoch, if replayed, produces a deviation shifted by $(M_1 - M_0)/\fref$
relative to the current reference. As $M$ grows monotonically
(Lemma~\ref{lem:mono}), this shift grows without bound, making the
replayed signal fall in a different cell or outside $\calD$ entirely.

\begin{figure*}[!htbp]
    \centering
    \includegraphics[width=\textwidth]{figures/panel_01.png}
    \caption{\textbf{Composition Inflation Theorem:
    $T(n,d) = d\,(d+1)^{n-1}$ counts distinguishable
    labeled timing trajectories and grows exponentially in both
    the number of events $n$ and the dimension $d$.}
    (\textbf{A}) $T(n,d)$ versus $n$ on a logarithmic axis for
    $d = 1,\ldots,6$ (coloured lines).
    Each series is linear on the log scale, confirming the exponential
    growth law; the slope $(d+1)$ increases with $d$, spanning seven
    decades over $n = 1$ to $11$.
    (\textbf{B}) Three-dimensional surface of $\log_{10} T(n,d)$ over
    the grid $n \in [1,10]$, $d \in [1,8]$.
    The surface rises steeply in both dimensions; the curvature in
    $d$ reflects the $(d+1)^{n-1}$ factor compounding with $d$.
    (\textbf{C}) Growth ratio $T(n,d)/T(n{-}1,d)$ versus $n$ for
    $d = 2,\ldots,6$.
    Dashed horizontal lines at $d{+}1$ confirm that the ratio is
    constant and equal to $d{+}1$ for all $n \geq 2$, providing
    direct numerical verification of the recurrence relation
    $T(n,d) = (d{+}1)\,T(n{-}1,d)$.
    (\textbf{D}) Direct enumeration versus closed-form formula for
    $d = 3$, $n = 1,\ldots,7$ (blue and red bars, log scale).
    The enumeration sums $d^k$ over all integer compositions of $n$
    into $k$ parts; the formula $T(n,3) = 3 \cdot 4^{n-1}$ matches
    exactly, with bars overlapping at every value of $n$.}
    \label{fig:composition_inflation}
\end{figure*}

%=======================================================================
\section{Security Analysis}
\label{sec:security}
%=======================================================================

\subsection{Threat Model}

We define an adversary $\mathcal{ADV}$ with the following
capabilities:
\begin{enumerate}[leftmargin=1.5em,topsep=2pt,itemsep=1pt]
  \item[\textbf{C1}] \emph{Content control}: $\mathcal{ADV}$ may
    choose arbitrary bit patterns for signal content.
  \item[\textbf{C2}] \emph{Timing control}: $\mathcal{ADV}$ may
    shift arrival times of legitimate pulses within $\pm \Delta_{\max}$.
  \item[\textbf{C3}] \emph{Injection}: $\mathcal{ADV}$ may inject
    additional pulses.
  \item[\textbf{C4}] \emph{Deletion}: $\mathcal{ADV}$ may suppress
    pulses.
  \item[\textbf{C5}] \emph{Replay}: $\mathcal{ADV}$ may retransmit
    recorded signals.
\end{enumerate}

This is a strong model: it subsumes Dolev-Yao for the temporal channel
and is strictly stronger than passive eavesdropping. We do not assume
cryptographic protection unless explicitly invoked.

\subsection{Payload Injection Impossibility}

\begin{theorem}[Structural Incorruptibility]
\label{thm:nopayload}
In a temporal programming system, capability \textbf{C1} (content
control) provides no leverage: the adversary cannot inject an action
not in the pre-compiled cell-action map $\mathcal{A}$.
\end{theorem}
\begin{proof}
In a temporal programming system, a received signal $e = (\chi, \DP)$
has two components: the channel identifier $\chi$ and the timing
deviation $\DP$. The channel identifier is determined by the physical
wire/frequency on which the pulse arrives (hardware-level), not by any
content. The timing deviation $\DP = \Tref(k) - \trec$ depends only
on arrival time, not amplitude, duration, or bit pattern.

Therefore, the adversary's content-control capability \textbf{C1}
affects neither $\chi$ nor $\DP$. Signal processing consists only of
the comparison $\DP \in \calC_i?$ for each $\calC_i \in \calP$, which
involves no parsing, no buffer operations, and no code interpretation.
The cell-action map $\mathcal{A}$ is compiled at deploy time and
stored in read-only memory (ROM). No runtime path leads from signal
content to $\mathcal{A}$ modification.

Consequently, the set of executable actions is always
$\{\mathcal{A}(\calC_i) : \calC_i \in \calP\}$, a finite set fixed
at compile time. $\mathcal{ADV}$ cannot expand this set.
\qed
\end{proof}

\begin{corollary}[Classical Injection Attacks Fail]
\label{cor:noinjection}
Buffer overflow attacks, SQL injection attacks, and command injection
attacks all require a payload parser that places adversary-supplied
data in a position where it is subsequently interpreted as code or
as a query \cite{saltzer1975protection}. In a temporal programming
system, no such parser exists. Therefore, these attack classes are
architecturally precluded, not merely mitigated.
\end{corollary}

\subsection{Replay Attack Failure}

\begin{theorem}[Replay Failure under Monotone $M$]
\label{thm:replay}
Let $\mathcal{ADV}$ record a legitimate trajectory $\tau =
(\DP_1^*, \ldots, \DP_n^*)$ at cycle count $M_0$. If $\mathcal{ADV}$
replays $\tau$ at any later time when $\Mcyc(t) = M_1 > M_0$, the
resulting trajectory $\tau' = (\DP_1', \ldots, \DP_n')$ satisfies
$\DP_j' = \DP_j^* + (M_1 - M_0)/\fref$ for each $j$.
In particular, $\tau'$ falls in the same cell as $\tau$ only if
$(M_1 - M_0)/\fref < \min_i \mathrm{width}(\calC_i)$.
\end{theorem}
\begin{proof}
At replay time, the reference oscillator has advanced to cycle $M_1$.
The $j$-th replayed pulse has the same physical arrival time
$t_{\mathrm{rec},j} = T_{\mathrm{ref},j}^* - \DP_j^*$ (since it is a physical replay).
But the current reference time for the $j$-th expected pulse is
$T_{\mathrm{ref},j}' = T_{\mathrm{ref},j}^* + (M_1 - M_0)/\fref$, since the oscillator
has advanced by $M_1 - M_0$ additional cycles.
Thus:
\begin{align}
  \DP_j' &= T_{\mathrm{ref},j}' - t_{\mathrm{rec},j} \notag \\
         &= \left[T_{\mathrm{ref},j}^* + \frac{M_1 - M_0}{\fref}\right]
            - \left[T_{\mathrm{ref},j}^* - \DP_j^*\right] \notag \\
         &= \DP_j^* + \frac{M_1 - M_0}{\fref}.
  \label{eq:replay-shift}
\end{align}
The shift $\delta = (M_1 - M_0)/\fref$ is constant across all events
in the replayed trajectory. If $\delta \geq \min_i \mathrm{width}(\calC_i)$,
at least one event exits its original cell. By Lemma~\ref{lem:mono},
$M_1 > M_0$ always holds for replay occurring after recording, so
$\delta > 0$.

As $M$ grows, $\delta \to \infty$, and with probability 1 over
uniformly distributed cell widths, the shifted trajectory exits
its original cell partition (or exits $\calD$ entirely), triggering
no action or an anomaly action. \qed
\end{proof}

\begin{remark}
The quantitative threshold $\delta < \min_i \mathrm{width}(\calC_i)$
provides a design equation: the minimum cell width should exceed the
maximum expected replay shift. For a 10~MHz crystal and cell width
100~ns, replay within $\lfloor 10^6 \cdot 100 \times 10^{-9} \rfloor = 1$
cycle succeeds. After 2 cycles ($\approx$200~ns), the shift exceeds
the cell width and replay fails. With a 10~MHz oscillator and 1~ms
minimum cell width, a replay attack must execute within 10,000~cycles
(1~ms), a very tight window.
\end{remark}

\begin{figure*}[!htbp]
    \centering
    \includegraphics[width=\textwidth]{figures/panel_03.png}
    \caption{\textbf{Monotone Time Coordinate and Structural Replay
    Resistance: $M(k)$ is non-decreasing by construction, and any
    replayed signal produces a $\Delta P$ distinguishable from the
    original.}
    (\textbf{A}) Monotone counter $M(k)$ versus signal index $k$
    (blue staircase with shaded area).
    Increments are drawn from $\{0,1,2\}$ at each step; the sequence
    is non-decreasing throughout, illustrating the constitutive
    asymmetry $t = M/f$ that forbids time reversal.
    (\textbf{B}) Timing deviation $\Delta P(k) = T_{\mathrm{ref}}(k)
    - t_{\mathrm{received}}(k)$ versus signal index.
    The signal fluctuates around zero with sub-cycle jitter; the
    shaded area tracks the sign of each deviation, providing the raw
    datum processed by timing cells.
    (\textbf{C}) Three-dimensional trajectory $(k,\,M_k,\,\Delta P_k)$
    in trajectory space.
    Red dots mark every tenth sample; the path is monotone in $M$
    regardless of $\Delta P$ sign, confirming that no trajectory can
    reverse its $M$-coordinate.
    (\textbf{D}) Scatter of replayed $\Delta P$ versus original $\Delta P$
    for 300 signal pairs.
    All replay values are shifted upward by a random delay
    $\delta \in [0.18,\,0.75]$~s, placing every point above the
    identity diagonal (dashed).
    A replayed signal is thus mapped to a different cell than the
    original, making the instruction structurally unexecutable.}
    \label{fig:monotone_replay}
\end{figure*}

\subsection{Bounded Attack Surface}

\begin{theorem}[Finite Enumerable Attack Surface]
\label{thm:surface}
The attack surface of a deployed temporal programming system is the
finite set $\Cells(\mathcal{A}) = \{\calC_1, \ldots, \calC_m\}$ of
pre-compiled cells. This set is:
\begin{enumerate}[leftmargin=1.5em,topsep=2pt,itemsep=1pt]
  \item[(i)] \emph{Finite}: $|\Cells(\mathcal{A})| = m < \infty$,
  \item[(ii)] \emph{Fixed at deploy time}: no runtime event can add
    or remove a cell (the registry is stored in ROM),
  \item[(iii)] \emph{Enumerable by the defender} before deployment.
\end{enumerate}
\end{theorem}
\begin{proof}
(i): The Tempus compiler produces a finite cell-action map $\mathcal{A}$
because the source program contains finitely many \textbf{cell}
declarations. (ii): Rule~\eqref{rule:assign} only looks up cells in
$\Gamma$ (the cell registry); no rule modifies $\Gamma$. $\Gamma$
is thus invariant under all reduction steps, consistent with
implementation in ROM. (iii): The compiler outputs $\Cells(\mathcal{A})$
as part of the build artifact, making it available for pre-deployment
security audit. \qed
\end{proof}

\begin{table}[ht]
\centering
\caption{Security comparison: temporal vs.\ content-based systems.}
\label{tab:security}
\small
\begin{tabular}{@{}lcc@{}}
\toprule
\textbf{Attack class} & \textbf{Content-based} & \textbf{Temporal} \\
\midrule
Buffer overflow     & Possible (requires parser)  & Impossible \\
SQL injection       & Possible (requires parser)  & Impossible \\
Command injection   & Possible (requires parser)  & Impossible \\
Replay              & Possible (no epoch check)   & Fails ($\delta$ shift) \\
Timing oracle       & Possible (side channel)     & By design only \\
Payload mutation    & Possible                    & Impossible \\
DoS (flooding)      & Possible                    & Rate-limited by $\fref$ \\
Attack surface      & Unbounded (all parse paths) & $|\Cells(\mathcal{A})|$ \\
\bottomrule
\end{tabular}
\end{table}

%=======================================================================
\section{Implementation Architecture}
\label{sec:impl}
%=======================================================================

\subsection{Hardware Layer}

The hardware layer provides the reference oscillator and pulse
reception circuitry.

\subsubsection{Oscillator Selection}

Three oscillator classes cover the relevant frequency and stability
range:

\begin{itemize}[leftmargin=1.5em,topsep=2pt,itemsep=1pt]
  \item \textbf{Crystal oscillators (XO/TCXO/OCXO):}
    Frequencies $1$--$100$~MHz, stability $\pm 10$~ppm (XO)
    to $\pm 0.01$~ppm (OCXO), cost \$0.50--\$200. Suitable for
    industrial IoT and embedded control.
  \item \textbf{MEMS oscillators:}
    Frequencies $1$--$125$~MHz, stability $\pm 50$~ppm,
    shock-resistant, cost \$1--\$10. Suitable for drones and
    mobile platforms.
  \item \textbf{Atomic clocks (Rb, Cs, optical):}
    Stability $<10^{-11}$/day, cost \$500--\$50\,000. Suitable for
    nuclear infrastructure and ground-truth nodes.
\end{itemize}

The minimum cell width $w_{\min}$ must satisfy:
\begin{equation}
  w_{\min} \;\geq\; \frac{2\epsilon}{\fref} + 2\delta_{\max}
  \label{eq:cellwidth}
\end{equation}
where $\epsilon$ is the fractional frequency stability and
$\delta_{\max}$ is the maximum propagation jitter.

\subsubsection{Pulse Reception and Timestamping}

A dedicated interrupt service routine (ISR) latches the oscillator
counter value at the rising edge of each received pulse.
The timestamp $\trec(k)$ is the latched counter divided by $\fref$.
This must complete within one oscillator period to avoid missed edges.
On a 10~MHz crystal, this is a 100~ns budget.

\subsection{Cell Registry}

The cell registry $\Gamma$ is a compile-time artifact stored in
read-only memory. Its structure is a balanced interval tree mapping
$\calD \to \calC_i$ for each channel.

\begin{algorithm}[ht]
\caption{Cell Registry Lookup}
\label{alg:lookup}
\begin{algorithmic}[1]
\Require $\DP$: timing deviation, $\chi$: channel, $\Gamma$: registry
\Ensure $\calC_i \in \calP_\chi$ such that $\DP \in \calC_i$, or $\bot$
\State $\mathit{tree} \gets \Gamma[\chi]$
\State $\mathit{node} \gets \mathit{tree}.\mathrm{root}$
\While{$\mathit{node} \neq \mathsf{nil}$}
  \If{$\DP < \mathit{node}.\mathit{lo}$}
    \State $\mathit{node} \gets \mathit{node}.\mathit{left}$
  \ElsIf{$\DP > \mathit{node}.\mathit{hi}$}
    \State $\mathit{node} \gets \mathit{node}.\mathit{right}$
  \Else
    \State \Return $\mathit{node}.\mathit{cell}$
  \EndIf
\EndWhile
\State \Return $\bot$
\end{algorithmic}
\end{algorithm}

Lookup is $O(\log m)$ for $m$ cells per channel. Since $m$ is small
(typically 4--32) and the tree is stored in ROM with cache-friendly
layout, lookup executes in 2--8 cycles on embedded processors.

\subsection{Trajectory Engine}

The trajectory engine accumulates events during the $\mathsf{COMPILE}$
phase and produces a cell assignment when the trajectory is complete.

\begin{algorithm}[ht]
\caption{Trajectory Engine}
\label{alg:traj}
\begin{algorithmic}[1]
\Require $n$: trajectory length, $\Gamma$: registry
\State $\tau \gets ()$; \quad $\sigma \gets \mathsf{IDLE}$
\Loop
  \State $\sigma \gets \mathsf{COMPILE}$
  \While{$|\tau| < n$}
    \State $e \gets \mathsf{WAIT\_EVENT}()$
      \Comment{blocking receive}
    \State $(\chi, \trec) \gets e$
    \State $k \gets \mathsf{CURRENT\_CYCLE}(\chi)$
    \State $\DP \gets \Tref(k) - \trec$
    \State $\calC \gets \mathsf{LOOKUP}(\DP, \chi, \Gamma)$
    \If{$\calC = \bot$}
      \State $\mathsf{LOG\_ANOMALY}(e)$; \textbf{goto} \textsc{Reset}
    \EndIf
    \State $\tau \gets \tau \cdot (\chi, \calC)$
  \EndWhile
  \State $\sigma \gets \mathsf{EXECUTE}$
  \State $\mathit{action} \gets \Gamma.\mathsf{COMPOSITE}(\tau)$
  \State $\mathsf{DISPATCH}(\mathit{action})$
  \label{alg:dispatch}
  \State \textsc{Reset}: $\tau \gets ()$; \quad $\sigma \gets \mathsf{IDLE}$
\EndLoop
\end{algorithmic}
\end{algorithm}

\subsection{Timing Scheduler and Phase State Machine}

\begin{figure}[ht]
\centering
\small
\begin{tabbing}
\hspace{3em}\=\kill
\texttt{IDLE}
  \>$\xrightarrow{\text{start}}$
  \texttt{COMPILE} \\[4pt]
\texttt{COMPILE}
  \>$\xrightarrow{|\tau|=n,\,\calC\neq\bot}$
  \texttt{EXECUTE} \\[4pt]
\texttt{COMPILE}
  \>$\xrightarrow{\calC=\bot}$
  \texttt{IDLE} (anomaly) \\[4pt]
\texttt{EXECUTE}
  \>$\xrightarrow{\phi=\varepsilon}$
  \texttt{IDLE}
\end{tabbing}
\caption{Phase state machine for the Tempus runtime.}
\label{fig:fsm}
\end{figure}

The phase state machine (Figure~\ref{fig:fsm}) implements
Theorem~\ref{thm:phase}: only one phase is active at any time,
and the $\mathsf{COMPILE} \leftrightarrow \mathsf{EXECUTE}$
transition is mediated by the completion of the trajectory.

\subsection{Dispatcher}

The dispatcher executes compiled actions. Since all actions are
pre-compiled (e.g., MISRA-C \cite{misrac2012} function calls or
DO-178C \cite{do178c} verified modules), dispatch is a table
lookup followed by a function call via a read-only function pointer
table. No runtime code generation, JIT compilation, or dynamic
loading occurs.

%=======================================================================
\section{Applications}
\label{sec:apps}
%=======================================================================

\subsection{Nuclear Plant Process Control}

Nuclear process control requires the highest levels of determinism
and tamper-resistance. Coolant flow, temperature, pressure, and
neutron flux must each be monitored continuously, with alarm latency
below 100~ms and zero tolerance for spurious actuation.

A temporal programming deployment maps each physical quantity to a
dedicated channel via an appropriate transducer
(Examples~\ref{ex:temp}, \ref{ex:press}). The cell partition for
each channel is designed using the formula \eqref{eq:cellwidth}
with OCXO oscillators ($\epsilon \approx 10^{-8}$) to achieve
cell widths of $\sim$1~ns---far narrower than physical measurement
noise, eliminating spurious cross-cell events.

The finite, enumerable attack surface (Theorem~\ref{thm:surface})
satisfies nuclear regulatory requirements for deterministic behavior.
No network stack is required: communication runs over dedicated
coaxial timing channels with no IP layer and no content parsing.

\subsection{Drone Swarm Autonomy}

Autonomous drone swarms require coordinated trajectory decisions
under adversarial radio environments. Current systems use encrypted
data-link protocols (MAVLink, DDS-RTPS), which remain vulnerable
to cryptanalytic attack and parsing errors in firmware.

A temporal programming approach uses the \textbf{ensemble}
declaration (Listing~\ref{lst:drone}) to synchronize swarm state
via shared timing pulses. Each drone measures $\DP$ relative to its
local MEMS oscillator, which is periodically synchronized to a
ground-truth atomic clock node (Theorem~\ref{thm:replay} provides
anti-replay protection for the synchronization pulses).

The $d=4$ channel, $n=3$ event trajectory supports
$T(3,4) = 4 \cdot 5^2 = 100$ distinguishable swarm states---adequate
for cooperative navigation. The entire swarm control program is
$<$4~KB of ROM, suitable for resource-constrained flight controllers.

\subsection{Industrial IoT Sensor Networks}

Industrial IoT networks (IIoT) often span thousands of sensor nodes
in hostile electromagnetic environments (foundries, chemical plants).
TinyOS-based \cite{levis2005tinyos} and similar systems face both
wireless interference and potential insider attacks on the data plane.

In a temporal programming deployment, sensor nodes transmit only
timing pulses over a carrier channel. No MAC-layer payload exists.
The cell registry maps chemical concentration ranges (via resonant
sensors) to alarm tiers. A node with a compromised firmware image
can at most cause mis-timing of its pulses by $\Delta_{\max}$;
it cannot inject arbitrary alarms (Theorem~\ref{thm:nopayload}).

Multi-domain composition (Proposition~\ref{prop:multidom}) allows a
single trajectory engine to process temperature, pressure, and
chemical concentration simultaneously over $d=3$ channels, producing
$T(n,3) = 3 \cdot 4^{n-1}$ distinguishable process states.

\begin{figure*}[!htbp]
    \centering
    \includegraphics[width=\textwidth]{figures/panel_05.png}
    \caption{\textbf{Multi-Domain Transduction: temperature, pressure,
    and light sensors each produce a scalar $\Delta P$ through
    domain-specific oscillator circuits, enabling composition of
    heterogeneous physical signals in a unified timing space.}
    (\textbf{A}) TCXO frequency deviation $\Delta P$ (in parts per
    million) versus ambient temperature over the industrial range
    $[-40,\,+85]$\,°C.
    The parabolic response $\Delta f/f \propto (T - T_0)^2$ maps
    the full temperature domain to a bounded $\Delta P$ interval,
    with the minimum at the turnover temperature $T_0 = 25$\,°C.
    (\textbf{B}) MEMS capacitive-to-timing delay $\Delta P$ (in
    nanoseconds) versus barometric pressure.
    The response is linear over $[0,\,1013]$\,hPa via the RC
    relation $\Delta\tau = R\,\Delta C(P)$, demonstrating monotone
    pressure-to-timing transduction.
    (\textbf{C}) Three-dimensional $\Delta P$ point cloud for 600
    simultaneous readings across temperature ($\Delta P_T$),
    pressure ($\Delta P_P$), and photodiode-RC ($\Delta P_L$) channels.
    Points are coloured by Euclidean distance from the origin;
    the cloud occupies a compact region confirming that all three
    dimensions are jointly bounded.
    (\textbf{D}) Density-normalised histograms of $\Delta P$ from
    four sensor types (TCXO, MEMS, RC-photodiode, piezo-vibration)
    overlaid on a common axis.
    Despite different physical origins, all distributions are
    centred near zero with characteristic widths, making them
    directly composable via the $D$-dimensional S-entropy framework.}
    \label{fig:multidomain}
\end{figure*}

%=======================================================================
\section{Evaluation}
\label{sec:eval}
%=======================================================================

\subsection{Computational Complexity}

\begin{table}[ht]
\centering
\caption{Per-event runtime complexity for key operations.}
\label{tab:complexity}
\small
\begin{tabular}{@{}lcc@{}}
\toprule
\textbf{Operation} & \textbf{Temporal} & \textbf{Content-based} \\
\midrule
Signal reception       & $O(1)$         & $O(L)$ parse \\
Cell lookup            & $O(\log m)$    & $O(1)$ (hash) \\
Trajectory accumulation& $O(1)$         & $O(L)$ buffer \\
Action dispatch        & $O(1)$         & $O(1)$ \\
Anomaly detection      & $O(\log m)$    & $O(L^2)$ introspection \\
Memory (state)         & $O(n \cdot d)$ & $O(n \cdot L)$ \\
Attack surface audit   & $O(m)$         & Undecidable (parser) \\
\bottomrule
\multicolumn{3}{l}{$L$ = message length, $m$ = cell count, $n$ = trajectory length.}
\end{tabular}
\end{table}

\subsection{Timing Overhead}

The timing overhead of a temporal programming node consists of:
\begin{enumerate}[leftmargin=1.5em,topsep=2pt,itemsep=1pt]
  \item \textbf{ISR latency}: 50--200~ns on ARM Cortex-M4 at 168~MHz.
  \item \textbf{Cell lookup}: 8--32 cycles ($\sim$50--200~ns at 168~MHz)
    for $m \leq 32$ cells.
  \item \textbf{Trajectory accumulation}: $O(1)$, $<$10 cycles.
  \item \textbf{Dispatch}: 2--5 cycles (function pointer call).
\end{enumerate}
Total per-event overhead: $< 1~\mu$s on commodity embedded hardware,
compared to $10$--$1000~\mu$s for TCP/IP stack processing of
equivalent data volumes \cite{levis2005tinyos,zhao2007ptides}.

\subsection{Security Overhead}

Security in temporal programming is achieved at zero marginal runtime
cost: the same cell lookup that classifies the event also enforces
the security policy. There is no separate cryptographic check, no
MAC verification, and no certificate validation. The \emph{structural}
security of Theorem~\ref{thm:nopayload} is free.

\subsection{Comparison with Related Systems}

\begin{table}[ht]
\centering
\caption{Feature comparison with related paradigms.}
\label{tab:comparison}
\small
\begin{tabular}{@{}lccccc@{}}
\toprule
\textbf{Feature} & \rotatebox{70}{\textbf{TTA}} &
  \rotatebox{70}{\textbf{PTIDES}} & \rotatebox{70}{\textbf{TinyOS}} &
  \rotatebox{70}{\textbf{Esterel}} & \rotatebox{70}{\textbf{Tempus}} \\
\midrule
No content parsing      & \texttimes & \texttimes & \texttimes & \texttimes & \checkmark \\
Payload-free channel    & \texttimes & \texttimes & \texttimes & \texttimes & \checkmark \\
Structural incorruptibility & \texttimes & \texttimes & \texttimes & \texttimes & \checkmark \\
Monotone replay defense & \texttimes & \texttimes & \texttimes & \texttimes & \checkmark \\
Formal cell semantics   & \texttimes & \texttimes & \texttimes & \texttimes & \checkmark \\
Deterministic timing    & \checkmark & \checkmark & \texttimes & \checkmark & \checkmark \\
Multi-domain sensing    & \checkmark & \checkmark & \checkmark & \texttimes & \checkmark \\
Synchronous semantics   & \checkmark & partial & \texttimes & \checkmark & \checkmark \\
Formal security proofs  & \texttimes & \texttimes & \texttimes & \texttimes & \checkmark \\
\bottomrule
\end{tabular}
\end{table}

\begin{figure*}[!htbp]
    \centering
    \includegraphics[width=\textwidth]{figures/panel_06.png}
    \caption{\textbf{Structural Incorruptibility: absence of a payload
    parser eliminates arbitrary code injection; the attack surface is
    a finite, pre-compiled cell-to-action map with zero novel-action
    capacity.}
    (\textbf{A}) Log-scale density histogram of 100\,000 adversarial
    payload sizes drawn from an exponential distribution
    with mean 500 bytes.
    Payloads span three decades; the log-scale axis reveals the
    heavy-tailed structure representative of real-world injection
    attempts.
    (\textbf{B}) Bar chart of input counts across five attack
    categories (noise, malformed, oversized, bad-tag, valid~$\Delta P$).
    The black legend entry ``novel actions $= 0$'' confirms that
    across all 100\,000 inputs, no action outside the pre-compiled
    cell-action map was executed.
    (\textbf{C}) Three-dimensional surface of detection latency
    (in clock cycles) over the joint attack parameter space
    (payload size $\times$ pattern complexity).
    The surface is flat at approximately one cycle across the entire
    domain, demonstrating $\mathcal{O}(1)$ rejection regardless of
    payload size or structural complexity.
    (\textbf{D}) Density-normalised histograms of $\Delta P$ measured
    during benign operation (blue) and under active adversarial
    injection (red).
    The two distributions are indistinguishable, confirming that
    payload content exerts no influence on timing behaviour:
    the attack channel and the signal channel are structurally
    disjoint.}
    \label{fig:incorruptibility}
\end{figure*}

%=======================================================================
\section{Discussion}
\label{sec:discussion}
%=======================================================================

\subsection{Limitations}

\textbf{Timing resolution.} Cell width is bounded below by oscillator
stability and propagation jitter. For commodity XO oscillators
($\epsilon = 10$~ppm), cells narrower than $\sim$10~ns are
unreliable. Atomic oscillators remove this bound at significant cost.

\textbf{Channel capacity.} Each channel carries exactly one bit of
information per event: which cell $\DP$ falls in. For $m$ cells,
the channel capacity is $\log_2 m$ bits per event by Shannon's
theorem \cite{shannon1948communication}. Rich data transmission
requires many events or many channels, making temporal programming
unsuitable for high-bandwidth data transfer.

\textbf{Clock drift.} Long-running systems require periodic
re-synchronization (NTP \cite{mills1991ntp} or PTP
\cite{ieee1588}) to prevent cells from drifting relative to the
physical phenomenon being measured. The synchronization period must
be shorter than the time for drift to exceed $w_{\min}/2$.

\textbf{Multi-hop networks.} Propagation delay accumulates over
network hops, adding to $\DP$. Each hop must either compensate for
its delay or absorb it within a cell's tolerance. PTIDES-style
timestamp forwarding \cite{zhao2007ptides} could be adapted.

\subsection{Relationship to Cyber-Physical Systems}

Lee \cite{lee2008cyber} argues that cyber-physical systems require
a tighter coupling between computational and physical time than
current software abstractions provide. Temporal programming takes
this argument to its logical conclusion: the physical quantity
\emph{is} the timing, and the program \emph{is} the cell partition.
No abstraction layer interposes between the physical phenomenon and
the computation.

\subsection{Future Work}

Immediate open problems include: (i) a type system for Tempus that
statically verifies cell coverage and disjointness; (ii) a compilation
target for FPGA-based cell lookup (sub-nanosecond lookup using LUT
cascades); (iii) formal verification of the Tempus runtime in a
proof assistant (Coq, Lean); (iv) extension to Byzantine fault
tolerance under capability \textbf{C4} (pulse deletion).

%=======================================================================
\section{Conclusion}
\label{sec:conclusion}
%=======================================================================

We have presented \emph{temporal programming}, a programming paradigm
in which every inter-node signal is a bare clock pulse and the program
datum is the timing deviation $\DP(k) = \Tref(k) - \trec(k)$.
Programs are finite collections of timing cells---measurable subsets
of $\DP$-space---that partition the observable space and map each
region to a pre-compiled action.

The core theoretical contribution is the Composition Inflation
Theorem: $T(n,d) = d\,(1+d)^{n-1}$ labeled trajectories are
distinguishable for $n$ events on $d$ channels, proved from first
principles via the stars-and-bars bijection and the binomial theorem.
This exponential growth in state space is achieved without transmitting
any content.

The central security result is Structural Incorruptibility: because no
payload parser exists, arbitrary code injection is architecturally
impossible. Replay attacks fail under monotone cycle count
(Theorem~\ref{thm:replay}), and the attack surface is a finite,
enumerable set of pre-compiled cells (Theorem~\ref{thm:surface}).
These properties hold by construction, not by mitigation.

The Tempus language provides a concrete realization of these ideas
with formal BNF syntax, small-step operational semantics, and a
provably phase-separated COMPILE/EXECUTE state machine. Applications
in nuclear control, drone autonomy, and industrial IoT demonstrate
the breadth of the paradigm.

Temporal programming does not replace content-rich protocols for
general-purpose computing. It offers a new foundation for the
growing class of safety-critical cyber-physical systems where
the need for structural immunity to attack is more important than
the ability to transmit arbitrary data.

%=======================================================================
% Acknowledgements
%=======================================================================
\section*{Acknowledgements}
The authors thank the Chitungwiza Institute for Distributed Systems
for supporting this research and the anonymous reviewers for detailed
comments on earlier drafts.

%=======================================================================
% References
%=======================================================================
\bibliographystyle{IEEEtran}
\bibliography{references}

\end{document}
